\apendice{Especificación de Requisitos}

\section{Introducción}

En este apartado se recoge la especificación de requisitos del sistema desarrollado. En él se describen los objetivos generales del proyecto (ver sección~\hyperref[req:objgen]{B.2}), así como los requisitos funcionales y no funcionales(ver sección~\hyperref[req:req]{B.3}) que definen el comportamiento esperado de la aplicación y las restricciones que debe cumplir durante su funcionamiento (ver sección~\hyperref[req:esp]{B.4}).

El propósito de esta especificación es servir como guía durante el desarrollo del sistema y como referencia para verificar que la aplicación implementa correctamente las funcionalidades previstas. Además, este documento permite definir de forma clara el alcance del proyecto y establecer una descripción detallada de las capacidades que ofrece la aplicación.

La organización de este apartado se ha realizado tomando como referencia las recomendaciones del estándar \file{IEEE 830-1998}~\cite{req:ieee} para la elaboración de documentos de especificación de requisitos software (\eng{Software Requirements Specification}, SRS). Dicho estándar establece una serie de características deseables para una correcta especificación de requisitos, indicando que esta debe ser completa, consistente, inequívoca, correcta, trazable, verificable y fácilmente modificable.

\section{Objetivos generales} \label{req:objgen}

Los objetivos generales que se buscan cumplir con este proyecto son:
\begin{itemize}
\tightlist
	\item Desarrollar un sistema basado en Generación Aumentada por Recuperación (\eng{RAG}) que permita realizar consultas sobre licitaciones del sector público, enriqueciendo las respuestas mediante información documental relevante.
	\item Automatizar la recopilación, procesamiento e indexación de documentos de licitaciones, de forma que puedan ser integrados eficientemente en una base de datos vectorial.
	\item Facilitar la consulta de la información de las licitaciones mediante una aplicación web intuitiva y accesible para usuarios autentificados.
	\item Mejorar la comprensión y análisis de la información administrativa y técnica contenida en las licitaciones.
	\item Proporcionar la trazabilidad de la respuesta del modelo, mostrando la información documental con la que se alimenta al \eng{LLM}. 
	\item Incorporar mecanismos automáticos de evaluación que permitan analizar la calidad de las respuestas generadas por el sistema \eng{RAG}.
\end{itemize}

\section{Catálogo de requisitos} \label{req:req}

A continuación se van a definir los requisitos del proyecto, tanto los funcionales como los no funcionales.

\textbf{Requisitos funcionales}

\begin{itemize}
\tightlist
	\item \req{RF-1} Gestión de documentos: el sistema debe permitir que el administrador gestione los documentos de licitaciones del sistema. 
	\begin{itemize}
	\tightlist
		\item \req{RF-1.1} Inserción de documentos: el sistema debe permitir cargar documentos de licitaciones con el objetivo de construir y actualizar la base de datos vectorial.
		\item \req{RF-1.2} Listado de documentos cargados: el sistema debe permitir ver los documentos cargados y sus metadatos.
		\item \req{RF-1.3} Borrado de documentos: el sistema debe permitir borrar los documentos cargados.
		\item \req{RF-1.4} Descarga de documentos: el sistema debe permitir descargar los documentos cargados.
	\end{itemize}
	\item \req{RF-2} Procesado de los documentos: el sistema debe procesar de manera automática los documentos introducidos para prepararlos para que el \eng{RAG} pueda usarlos.
	\begin{itemize}
	\tightlist
		\item \req{RF-2.1} Extracción de información: el sistema debe extraer automáticamente el contenido textual de los documentos introducidos, ignorando páginas o documentos sin texto extraíble.
		\item \req{RF-2.2} Conversión del contenido: el sistema debe transformar el texto extraído a un formato homogéneo, concretamente a \eng{markdown} para que sea más fácil su procesamiento.
		\begin{itemize}
		\tightlist
			\item \req{RF-2.2.1} Conservar estructura básica: el sistema debe conservar la estructura básica del texto siempre que sea posible (secciones, títulos, saltos de línea, párrafos, etc).
		\end{itemize}
		\item \req{RF-2.3} Limpieza y normalización del texto: el sistema debe limpiar y normalizar el texto que se extrae de los documentos.
		\begin{itemize}
		\tightlist
			\item \req{RF-2.3.1} Eliminar caracteres: el sistema debe eliminar los caracteres irrelevantes o problemáticos.
			\item \req{RF-2.3.2} Normalización: el sistema debe normalizar los espacios, saltos de línea y formatos inconsistentes.
		\end{itemize}
	\end{itemize}
	\item \req{RF-3} Segmentación en \eng{chunks}: el sistema debe dividir los documentos ya procesados en fragmentos de texto (\eng{chunks}) adecuados para el modelo de \eng{embeddings}.
	\begin{itemize}
	\tightlist
		\item \req{RF-3.1} Tamaño: el sistema debe segmentar el texto respetando un máximo de tamaño medido en número de \eng{tokens} para evitar errores.
		\item \req{RF-3.2} Solapamiento: el sistema debe aplicar solapamiento entre \eng{chunks} para preservar el contexto.
		\item \req{RF-3.3} Contexto: cada \eng{chunk} debe tener una referencia al documento original y a su posición dentro de él.
	\end{itemize}
	\item \req{RF-4} Generación de \eng{embeddings}: el sistema debe generar representaciones vectoriales (\eng{embeddings}) para cada \eng{chunk}.
	\item \req{RF-5} Persistencia en base de datos vectorial: el sistema debe almacenar los \eng{chunks} y sus \eng{embeddings} en una base de datos vectorial.
	\begin{itemize}
	\tightlist
		\item \req{RF-5.1} Almacenar \eng{chunks}: el sistema debe almacenar el contenido textual de cada \eng{chunk}.
		\item \req{RF-5.2} Almacenar \eng{embeddings}: el sistema debe almacenar el \eng{embedding} asociado a cada \eng{chunk}.
		\item \req{RF-5.3} Almacenar metadatos: el sistema debe almacenar metadatos relevantes como el nombre del archivo, el título y el segmento incluido en el \eng{chunk}.
		\item \req{RF-5.4} Actualización: el sistema debe permitir la actualización o recreación de la colección vectorial en el caso de incluir documentos nuevos.
		\item \req{RF-5.5} Eliminar: el sistema debe permitir borrar los \eng{chunks} y \eng{embeddings} de la base de datos vectorial.
	\end{itemize}
	\item \req{RF-6} Recuperación de información por similitud: el sistema debe recuperar información relevante a partir de las consultas de los usuarios.
	\begin{itemize}
	\tightlist
		\item \req{RF-6.1} Generación de \eng{embedding}: el sistema debe generar el \eng{embedding} de la consulta del usuario.
		\item \req{RF-6.2} Búsqueda por similitud: el sistema debe realizar una búsqueda por similitud en la base de datos vectorial.
		\item \req{RF-6.3} Filtrado por metadatos: el sistema debe permitir filtrar al recuperación utilizando los metadatos asociados a los documentos.
	\end{itemize}
	\item \req{RF-7} Construcción del \eng{prompt} con contexto: el sistema debe construir un \eng{prompt}, para alimentar al \eng{LLM}, que incluya los \eng{chunks} recuperados, la pregunta del usuario e instrucciones de uso del contexto.
	\item \req{RF-8} Generación de respuestas mediante un \eng{LLM}: el sistema debe enviar el \eng{prompt} generado a un modelo \eng{LLM} y mostrar la respuesta de dicho modelo al usuario.
	\begin{itemize}
	\tightlist
		\item \req{RF-8.1} Selección de modelo: el sistema debe permitir utilizar distintos modelos \eng{LLM}.
		\item \req{RF-8.2} Control de alucinaciones: el sistema debe indicar cuando no existe suficiente información relevante para responder. 
	\end{itemize}
	\item \req{RF-9} Trazabilidad de respuesta: el sistema debe devolver junto con la respuesta los metadatos del documento origen y los \eng{chunks} empleados para la generación de la respuesta, así como su similitud. 
	\item \req{RF-10} Gestión de usuarios: el sistema debe permitir la gestión de usuarios a través de la aplicación web.
	\begin{itemize}
	\tightlist
		\item \req{RF-10.1} Registros y autentificación: el sistema debe permitir el registro y autentificación de los usuarios.
		\item \req{RF-10.2} Sesiones: el sistema debe gestionar las sesiones de los usuarios de forma segura.
		\item \req{RF-10.3} Roles: el sistema debe permitir asignar a los usuarios el rol de administradores o usuarios normales.
		\item \req{RF-10.4} Administradores: los administradores deben poder añadir usuarios, eliminarlos y cambiar su rol.
		\item \req{RF-10.5} Edición de datos de usuario: los usuarios deben poder editar sus datos.
		\item \req{RF-10.6} Recuperación de contraseña: los usuarios deben poder cambiar su contraseña utilizando su correo.
	\end{itemize}
	\item \req{RF-11} Interfaz de consultas: los usuarios autentificados podrán realizar consultas y ver las respuestas del sistema.
	\begin{itemize}
	\tightlist
		\item \req{RF-11.1} Historial de consultas: los usuarios podrán ver su historial de consultas.
		\item \req{RF-11.2} Borrar consultas: los usuarios podrán borrar consultas del historial.
		\item \req{RF-11.3} Consultas guiadas: el sistema debe permitir realizar consultas guiadas especializadas mediante \eng{prompts} predefinidos. Debe incluir resúmenes, importes, solvencia, garantías, criterios de adjudicación y otras preguntas típicas de las licitaciones. 
	\end{itemize}
	\item \req{RF-12} Importación automática de licitaciones: el sistema debe permitir al administrador importar de manera automática las licitaciones mediante \eng{web scraping}.
	\begin{itemize}
	\tightlist
		\item \req{RF-12.1} Descarga y almacenamiento: el sistema debe descargar y almacenar los documentos obtenidos, así como sus metadatos.
		\item \req{RF-12.2} Gestión de errores y duplicados: el sistema debe continuar aunque la descarga de algún documento falle. Debe registrar los errores y evitar duplicados.
	\end{itemize}
	\item \req{RF-13} Mostrar estadísticas: el sistema debe permitir ver estadísticas y gráficos de uso y de comparación de modelos.
	\item \req{RF-14} Evaluación del sistema \eng{RAG}: el sistema debe permitir evaluar automáticamente la calidad de las respuestas generadas. 
	\begin{itemize}
	\tightlist
		\item \req{RF-14.1} Generación de \eng{datasets}: el sistema debe permitir generar \eng{dtasets} de evaluación automáticos.
	\end{itemize}
\end{itemize}

\textbf{Requisitos no funcionales}
\begin{itemize}
\tightlist
	\item \req{RNF-1} Reproducibilidad: el proyecto debe ejecutarse en un entorno reproducible.
	\item \req{RNF-2} Rendimiento: el sistema debe ser capaz de procesar los documentos sin intervención manual. Además, el sistema tiene que responder en un tiempo razonable.
	\item \req{RNF-3} Robustez ante errores: el sistema debe continuar ejecutándose ante fallos esperables en datos, red o servicios externos.
	\begin{itemize}
	\tightlist
		\item \req{RNF-3.1} Documentos corruptos: el sistema debe saltarse los documentos corruptos o sin texto extraíble y continuar con el resto de documentos.
		\item \req{RNF-3.2} \eng{Timeouts}: el sistema debe manejar los \eng{timeouts} del \eng{LLM} devolviendo mensajes controlados al usuario.
	\end{itemize}
	\item \req{RNF-4} Trazabilidad técnica: el sistema debe realizar \eng{logs} de carga de modelo, lectura de los PDF, \eng{chunking}, \eng{embeddings}, guardado y consulta.
	\item \req{RNF-5} Seguridad: el sistema debe gestionar de forma segura la información sensible de los usuarios y el acceso a la aplicación.
	\begin{itemize}
	\tightlist
		\item \req{RNF-5.1} Contraseñas: las contraseñas deben almacenarse encriptadas.
		\item \req{RNF-5.2} Control de acceso por roles: las funciones administrativas deben estar protegidas por rol para que solo las pueda realizar el administrador.
		\item \req{RNF-5.3} Protección de secretos: las claves sensibles no deben estar públicas en el repositorio.
	\end{itemize}
	\item \req{RNF-6} Usabilidad: la interfaz debe permitir a los usuarios realizar las consultas y ver el historial de manera sencilla. Además, los textos deben ser legibles.
	\item \req{RNF-7} Calidad del \eng{software}: debe existir un conjunto de pruebas que asegure el correcto funcionamiento del programa. 
	\item \req{RNF-8} Disponibilidad: la web debe garantizar un nivel de disponibilidad adecuado que permita a los usuarios acceder a la aplicación y realizar consultas.
	\item \req{RNF-9} Recuperación y tolerancia a fallos: el sistema debe disponer de mecanismos que permitan la recuperación de información y la restauración del servicio tras fallos, garantizando la integridad de los datos.
	\item \req{RNF-10} Evitar alucinaciones: el sistema no debe responder cuando no haya contexto relevante suficiente.
	\item \req{RNF-11} Calidad de las respuestas del \eng{LLM}: el sistema debe llevar a cabo un control de calidad de las respuestas que da el \eng{LLM} a las consultas del usuario.
	\item \req{RNF-12} Escalabilidad: el sistema debe permitir incorporar nuevos documentos sin necesidad de reconstruir la aplicación.
	\item \req{RNF-13} Modularidad: el sistema debe permitir sustituir componentes como el \eng{LLM}, el modelo de \eng{embeddings} o la base de datos vectorial con cambios mínimos.
	\item \req{RNF-14} Portabilidad: la aplicación debe poder desplegarse en distintos entornos utilizando contenedores \file{Docker}.
	\item \req{RNF-15} Mantenibilidad: el sistema debe seguir una arquitectura modular y debe estar correctamente documentado para que se facilite su mantenimiento.
	\item \req{RNF-16} Internacionalización: el sistema debe traducirse de manera automática al inglés.
	\item \req{RNF-17} Compatibilidad \eng{responsive}: la interfaz debe adaptarse a distintos tamaños de pantalla.
	\item \req{RFN-18} Obsevabilidad: el sistema debe proporcionar \eng{logs} y métricas para monitorizar el estado del procesamiento y las consultas.
\end{itemize}

\section{Especificación de requisitos} \label{req:esp}

Los actores del sistema son:
\begin{itemize}
\tightlist
	\item \textbf{Usuario anónimo}: usuario que no ha iniciado sesión (solo puede registrarse y autentificarse)
	\item \textbf{Usuario autentificado}: usuario que puede acceder a la aplicación sus acciones en al aplicación cambiarán dependiendo de su rol. 
	\begin{itemize}
	\tightlist
	 	\item \textbf{Usuario normal}: usuario autentificado que puede consultar el \eng{RAG} y ver su historial
	 	\item \textbf{Usuario administrador}: usuario autentificado con permisos de gestión de usuarios y de carga y actualización documental
	 \end{itemize}
	 \item \textbf{Sistema}: es el encargado del procesamiento, \eng{chunking}, \eng{embeddings}, guardado y consulta. 
\end{itemize}

Los casos de uso definidos para el proyecto son:
\begin{itemize}
\tightlist
	\item \textbf{CU-1} Registro de usuarios (ver tabla~\ref{req-01}).
	\item \textbf{CU-2} Iniciar sesión (ver tabla~\ref{req-02}).
	\item \textbf{CU-3} Cerrar sesión (ver tabla~\ref{req-03}).
	\item \textbf{CU-4} Recuperar contraseña (ver tabla~\ref{req-04}).
	\item \textbf{CU-5} Realizar consulta al \eng{RAG} (ver tabla~\ref{req-05}).
	\item \textbf{CU-6} Consultar historial de consultas (ver tabla~\ref{req-06}).
	\item \textbf{CU-7} Borrar consultas del historial (ver tabla~\ref{req-07}).
	\item \textbf{CU-8} Gestión de documentos (ver tabla~\ref{req-08}).
	\item \textbf{CU-9} Cargar documentos de licitaciones (ver tabla~\ref{req-09}).
	\item \textbf{CU-10} Importar licitaciones automáticamente (ver tabla~\ref{req-10}).
	\item \textbf{CU-11} Convertir los documentos a \eng{Markdown} (ver tabla~\ref{req-11}).
	\item \textbf{CU-12} Indexar documento (ver tabla~\ref{req-12}).
	\item \textbf{CU-13} Listar documentos (ver tabla~\ref{req-13}).
	\item \textbf{CU-14} Consultar estado del documento (ver tabla~\ref{req-14}).
	\item \textbf{CU-15} Eliminar documento (ver tabla~\ref{req-15}).
	\item \textbf{CU-16} Descargar documento (ver tabla~\ref{req-16}).
	\item \textbf{CU-17} Gestión de usuario (ver tabla~\ref{req-17}).
	\item \textbf{CU-18} Crear usuario (ver tabla~\ref{req-18}).
	\item \textbf{CU-19} Borrar usuario (ver tabla~\ref{req-19}).
	\item \textbf{CU-20} Modificar rol del usuario (ver tabla~\ref{req-20}).
	\item \textbf{CU-21} Editar usuario (ver tabla~\ref{req-21}).
	\item \textbf{CU-22} Realizar consulta guiada (ver tabla~\ref{req-22}).
	\item \textbf{CU-23} Cambiar idioma (ver tabla~\ref{req-23}).
	\item \textbf{CU-24} Cambiar tema visual (ver tabla~\ref{req-24}).
	\item \textbf{CU-25} Visualizar estadísticas (ver tabla~\ref{req-25}).
	\item \textbf{CU-26} Ejecutar evaluación \eng{RAG} (ver tabla~\ref{req-26}).
\end{itemize}
El diagrama de los casos de uso se muestra en la imagen ~\ref{fig:Diagramas/CasosUso}.

\imagendiagram{Diagramas/CasosUso}{Diagrama general de casos de uso.}{1.2}


\begin{table}[p]
	\centering
	\begin{tabularx}{\linewidth}{ p{0.21\columnwidth} p{0.71\columnwidth} }
		\toprule
		\textbf{\textbf{CU-1}}    & \textbf{Registro de usuarios}\\
		\toprule
		\textbf{Versión}              & 1.0    \\
		\textbf{Autor}                & Lydia Blanco Ruiz \\
		\textbf{Requisitos asociados} & \refreq{RF-10.1} \\
		\textbf{Descripción}          & Permite crear una cuenta en la aplicación web. \\
		\textbf{Precondición}         & Que el usuario no esté autentificado. \\
		\textbf{Acciones}             &
		\begin{enumerate}
			\def\labelenumi{\arabic{enumi}.}
			\tightlist
			\item El usuario accede a <<Registro>>.
			\item El usuario introduce el nombre, \eng{email} y contraseña.
			\item El sistema valida el formato y que el \eng{email} no esté duplicado.
			\item El sistema almacena el usuario con la contraseña encriptada.
			\item El sistema confirma el registro.
			\item El sistema crea la sesión y redirige a la página de consultas.
		\end{enumerate}\\
		\textbf{Postcondición}        & Cuenta creada y usuario autentificado \\
		\textbf{Excepciones}          & 
		\begin{enumerate}
			\def\labelenumi{\arabic{enumi}.}
			\tightlist
			\item Tiene algún campo vacío $\Rightarrow$ Mostrar error de campo obligatorio.
			\item El \eng{email} ya existe $\Rightarrow$ Mostrar error de \eng{email} duplicado.
			\item El \eng{email} no cumple la validación $\Rightarrow$ Mostrar error de \eng{email} inválido.
			\item La contraseña no cumple la política $\Rightarrow$ Mostrar error de contraseña poco segura.
			\item Las contraseñas no coinciden $\Rightarrow$ Mostrar error de contraseñas distintas.
			\item El nombre tiene menos de 2 caracteres $\Rightarrow$ Mostrar error de nombre demasiado corto.
		\end{enumerate}\\
		\textbf{Importancia}          & Alta.  \\
		\bottomrule
	\end{tabularx}
	\caption{\textbf{CU-1} Registro de usuarios.}
	\label{req-01}
\end{table}

\begin{table}[p]
	\centering
	\begin{tabularx}{\linewidth}{ p{0.21\columnwidth} p{0.71\columnwidth} }
		\toprule
		\textbf{\textbf{CU-2}}    & \textbf{Iniciar sesión}\\
		\toprule
		\textbf{Versión}              & 1.0    \\
		\textbf{Autor}                & Lydia Blanco Ruiz \\
		\textbf{Requisitos asociados} & \refreq{RF-10.1}, \refreq{RF-10.2} \\
		\textbf{Descripción}          & Autentifica al usuario y abre una sesión segura. \\
		\textbf{Precondición}         & Que el usuario esté registrado y no esté autentificado. \\
		\textbf{Acciones}             &
		\begin{enumerate}
			\def\labelenumi{\arabic{enumi}.}
			\tightlist
			\item El usuario accede a <<Iniciar sesión>>.
			\item El usuario introduce el \eng{email} y la contraseña.
			\item El sistema verifica los credenciales.
			\item El sistema crea la sesión y redirige a la página de consultas.
		\end{enumerate}\\
		\textbf{Postcondición}        & Sesión activa para el usuario. \\
		\textbf{Excepciones}          & 
		\begin{enumerate}
			\def\labelenumi{\arabic{enumi}.}
			\tightlist
			\item Tiene algún campo vacío $\Rightarrow$ Mostrar error de campo obligatorio.
			\item El \eng{email} o la contraseña no coinciden con ningún usuario registrado $\Rightarrow$ Mostrar error de usuario o contraseña incorrectos.
			\item El \eng{email} no cumple la validación $\Rightarrow$ Mostrar error de \eng{email} inválido.
			\item La contraseña no cumple la política $\Rightarrow$ Mostrar error de contraseña poco segura.		
		\end{enumerate}\\
		\textbf{Importancia}          & Alta.  \\
		\bottomrule
	\end{tabularx}
	\caption{\textbf{CU-2} Iniciar sesión.}
	\label{req-02}
\end{table}

\begin{table}[p]
	\centering
	\begin{tabularx}{\linewidth}{ p{0.21\columnwidth} p{0.71\columnwidth} }
		\toprule
		\textbf{\textbf{CU-3}}    & \textbf{Cerrar sesión}\\
		\toprule
		\textbf{Versión}              & 1.0    \\
		\textbf{Autor}                & Lydia Blanco Ruiz \\
		\textbf{Requisitos asociados} & \refreq{RF-10.2} \\
		\textbf{Descripción}          & Finaliza la sesión del usuario. \\
		\textbf{Precondición}         & Que el usuario esté autentificado. \\
		\textbf{Acciones}             &
		\begin{enumerate}
			\def\labelenumi{\arabic{enumi}.}
			\tightlist
			\item El usuario pulsa el botón <<Cerrar sesión>>.
			\item El sistema invalida la sesión.
			\item El sistema redirige a la página de inicio.
		\end{enumerate}\\
		\textbf{Postcondición}        & Sesión finalizada. \\
		\textbf{Excepciones}          & 
		\begin{enumerate}
			\def\labelenumi{\arabic{enumi}.}
			\tightlist
			\item Error al invalidar sesión $\Rightarrow$ Mostrar mensaje de error al cerrar sesión. 	
		\end{enumerate}\\
		\textbf{Importancia}          & Media.  \\		
		\bottomrule
	\end{tabularx}
	\caption{\textbf{CU-3} Cerrar sesión.}
	\label{req-03}
\end{table}

\begin{table}[p]
	\centering
	\begin{tabularx}{\linewidth}{ p{0.21\columnwidth} p{0.71\columnwidth} }
		\toprule
		\textbf{\textbf{CU-4}}    & \textbf{Recuperar contraseña}\\
		\toprule
		\textbf{Versión}              & 1.0    \\
		\textbf{Autor}                & Lydia Blanco Ruiz \\
		\textbf{Requisitos asociados} & \refreq{RF-10}, \refreq{RF-10.1}, \refreq{RF-10.6} \\
		\textbf{Descripción}          & Permite a los usuarios cambiar su contraseña antes de iniciar sesión. \\
		\textbf{Precondición}         & Que el usuario esté registrado en el sistema. \\
		\textbf{Acciones}             &
		\begin{enumerate}
			\def\labelenumi{\arabic{enumi}.}
			\tightlist
			\item El usuario accede a <<Iniciar sesión>>.
			\item El usuario pulsa el enlace <<¿No recuerdas la contraseña?>>.
			\item El sistema envía un correo al email del usuario con el enlace a la página de recuperación de contraseña.
			\item El usuario accede a la página de recuperar contraseña.
			\item El usuario introduce la nueva contraseña y la confirmación de la contraseña.
			\item El sistema guarda la nueva contraseña encriptada.
		\end{enumerate}\\
		\textbf{Postcondición}        & Contraseña cambiada para el usuario. \\
		\textbf{Excepciones}          & 
		\begin{enumerate}
			\def\labelenumi{\arabic{enumi}.}
			\tightlist
			\item Error al enviar el correo $\Rightarrow$ Mostrar error al enviar el correo.
			\item La contraseña no cumple la política $\Rightarrow$ Mostrar error de contraseña poco segura.
			\item Las contraseñas no coinciden $\Rightarrow$ Mostrar error de contraseñas distintas.
		\end{enumerate}\\
		\textbf{Importancia}          & Media.  \\
		\bottomrule
	\end{tabularx}
	\caption{\textbf{CU-4} Recuperar contraseña.}
	\label{req-04}
\end{table}

\begin{table}[p]
	\centering
	\begin{tabularx}{\linewidth}{ p{0.21\columnwidth} p{0.71\columnwidth} }
		\toprule
		\textbf{\textbf{CU-5}}    & \textbf{Realizar consulta al \eng{RAG}}\\
		\toprule
		\textbf{Versión}              & 1.0    \\
		\textbf{Autor}                & Lydia Blanco Ruiz \\
		\textbf{Requisitos asociados} & \refreq{RF-6}, \refreq{RF-7}, \refreq{RF-8}, \refreq{RF-9}, \refreq{RF-11}, \refreq{RNF-3}, \refreq{RNF-3.2}, \refreq{RNF-10}, \refreq{RNF-11} \\
		\textbf{Descripción}          & El usuario hace una pregunta, el sistema recupera contexto por similitud, construye el \eng{prompt} y genera respuesta que incluya las fuentes. \\
		\textbf{Precondición}         & Que el usuario esté autentificado y la base de datos vectorial esté disponible. \\
		\textbf{Acciones}             &
		\begin{enumerate}
			\def\labelenumi{\arabic{enumi}.}
			\tightlist
			\item El usuario escribe y envía una pregunta.
			\item El sistema genera \eng{embeddings} de la consulta.
			\item El sistema busca por similitud en la base de datos vectorial.
			\item El sistema construye un \eng{prompt} con los \eng{chunks} y las reglas de uso del contexto.
			\item El sistema llama al \eng{LLM} y genera respuestas.
			\item El sistema adjunta la trazabilidad de la respuesta, la cual incluye los documentos, metadatos y \eng{chunks} utilizados.
			\item El sistema guarda la consulta en el historial del consultas del usuario.
		\end{enumerate}\\
		\textbf{Postcondición}        & Respuesta mostrada y consulta registrada. \\
		\textbf{Excepciones}          & 
		\begin{enumerate}
			\def\labelenumi{\arabic{enumi}.}
			\tightlist
			\item Cuando no hay contexto suficiente $\Rightarrow$ El sistema no responde a la pregunta (para evitar alucinaciones) y muestra un aviso de que falta contexto.
			\item \eng{Timeout} del \eng{LLM} $\Rightarrow$ Mostrar un mensaje controlado. 	
		\end{enumerate}\\
		\textbf{Importancia}          & Alta.  \\
		\bottomrule
	\end{tabularx}
	\caption{\textbf{CU-5} Realizar consulta al \eng{RAG}.}
	\label{req-05}
\end{table}

\begin{table}[p]
	\centering
	\begin{tabularx}{\linewidth}{ p{0.21\columnwidth} p{0.71\columnwidth} }
		\toprule
		\textbf{\textbf{CU-6}}    & \textbf{Consultar historial de consultas}\\
		\toprule
		\textbf{Versión}              & 1.0    \\
		\textbf{Autor}                & Lydia Blanco Ruiz \\
		\textbf{Requisitos asociados} & \refreq{RF-11.1} \\
		\textbf{Descripción}          & Permite al usuario ver su lista de preguntas y respuestas anteriores. \\
		\textbf{Precondición}         & Que el usuario esté autentificado y existan consultas guardadas. \\
		\textbf{Acciones}             &
		\begin{enumerate}
			\def\labelenumi{\arabic{enumi}.}
			\tightlist
			\item El usuario accede al <<Historial>>.
			\item El sistema lista las consultas del usuario.
			\item El usuario puede abrir una consulta para ver la respuesta completa y la trazabilidad de dicha respuesta.
		\end{enumerate}\\
		\textbf{Postcondición}        & Historial visualizado. \\
		\textbf{Excepciones}          & 
		\begin{enumerate}
			\def\labelenumi{\arabic{enumi}.}
			\tightlist
			\item El usuario no ha realizado consultas $\Rightarrow$ Mostrar mensaje de que aun no hay consultas. 	
		\end{enumerate}\\
		\textbf{Importancia}          & Alta.  \\
		\bottomrule
	\end{tabularx}
	\caption{\textbf{CU-6} Consultar historial de consultas.}
	\label{req-06}
\end{table}

\begin{table}[p]
	\centering
	\begin{tabularx}{\linewidth}{ p{0.21\columnwidth} p{0.71\columnwidth} }
		\toprule
		\textbf{\textbf{CU-7}}    & \textbf{Borrar consultas del historial}\\
		\toprule
		\textbf{Versión}              & 1.0    \\
		\textbf{Autor}                & Lydia Blanco Ruiz \\
		\textbf{Requisitos asociados} & \refreq{RF-11.2} \\
		\textbf{Descripción}          & Permite al usuario borrar consultas del historial. \\
		\textbf{Precondición}         & Que el usuario esté autentificado y existan consultas guardadas. \\
		\textbf{Acciones}             &
		\begin{enumerate}
			\def\labelenumi{\arabic{enumi}.}
			\tightlist
			\item El usuario accede al <<Historial>>.
			\item El sistema lista las consultas del usuario.
			\item El usuario pulsa el botón <<Borrar>>.
			\item El sistema elimina la consulta seleccionada de la base de datos.
			\item El sistema actualiza la vista del historial.
		\end{enumerate}\\
		\textbf{Postcondición}        & La consulta seleccionada se elimina del historial del usuario. \\
		\textbf{Excepciones}          & 
		\begin{enumerate}
			\def\labelenumi{\arabic{enumi}.}
			\tightlist
			\item Error en la base de datos al intentar borrar la consulta $\Rightarrow$ Mostrar mensaje de error controlado.
			\item La consulta seleccionada no existe o ya ha sido borrada $\Rightarrow$ Mostrar aviso al usuario.
		\end{enumerate}\\
		\textbf{Importancia}          & Alta.  \\
		\bottomrule
	\end{tabularx}
	\caption{\textbf{CU-7} Borrar consultas del historial.}
	\label{req-07}
\end{table}

\begin{table}[p]
	\centering
	\begin{tabularx}{\linewidth}{ p{0.21\columnwidth} p{0.71\columnwidth} }
		\toprule
		\textbf{\textbf{CU-8}}    & \textbf{Gestión de documentos}\\
		\toprule
		\textbf{Versión}              & 1.0    \\
		\textbf{Autor}                & Lydia Blanco Ruiz \\
		\textbf{Requisitos asociados} & \refreq{RF-1}, \refreq{RF-1.1}, \refreq{RF-1.2}, \refreq{RF-1.3}, \refreq{RF-2}, \refreq{RF-2.1}, \refreq{RF-2.2}, \refreq{RF-2.2.1}, \refreq{RF-2.3}, \refreq{RF-2.3.1}, \refreq{RF-2.3.2}, \refreq{RF-3}, \refreq{RF-4}, \refreq{RF-5}, \refreq{RF-5.1}, \refreq{RF-5.2}, \refreq{RF-5.3}, \refreq{RF-5.4}, \refreq{RF-5.5}, \refreq{RNF-3.1} \\
		\textbf{Descripción}          & El administrador debe poder gestionar a los documentos de la aplicación. \\
		\textbf{Precondición}         & Hay un administrador autentificado. \\
		\textbf{Acciones}             &
		\begin{enumerate}
			\def\labelenumi{\arabic{enumi}.}
			\tightlist
			\item El administrador accede a la página de <<Iniciar sesión>>.
			\item El administrador introduce su \eng{email} y contraseña.
			\item El sistema le autentifica como administrador.
			\item El administrador accede a <<Administrar documentos>>.
		\end{enumerate}\\
		\textbf{Postcondición}        & Posibilidad de cargar documentos nuevos (de forma manual o por \eng{web scraping}), indexar, borrar o descargar los documentos ya cargados . \\
		\textbf{Excepciones}          & 
		\begin{enumerate}
			\def\labelenumi{\arabic{enumi}.}
			\tightlist
			\item Intento de acceso sin ser administrador $\Rightarrow$ Se deniega.
			\item Si no hay documentos cargados en la aplicación $\Rightarrow$ Muestra un mensaje.
		\end{enumerate}\\
		\textbf{Importancia}          & Alta.  \\
		\bottomrule
	\end{tabularx}
	\caption{\textbf{CU-8} Gestión de documentos.}
	\label{req-08}
\end{table}

\begin{table}[p]
	\centering
	\begin{tabularx}{\linewidth}{ p{0.21\columnwidth} p{0.71\columnwidth} }
		\toprule
		\textbf{\textbf{CU-9}}    & \textbf{Cargar documentos de licitaciones}\\
		\toprule
		\textbf{Versión}              & 1.0    \\
		\textbf{Autor}                & Lydia Blanco Ruiz \\
		\textbf{Requisitos asociados} & \refreq{RF-1}, \refreq{RF-10.3}, \refreq{RF-10.4}, \refreq{RNF-3}, \refreq{RNF-3.1} \\
		\textbf{Descripción}          & El administrador sube documentos para alimentar o ampliar la base de datos vectorial. \\
		\textbf{Precondición}         & Que el administrador esté autentificado. \\
		\textbf{Acciones}             &
		\begin{enumerate}
			\def\labelenumi{\arabic{enumi}.}
			\tightlist
			\item El administrador pulsa el botón <<Cargar>>.
			\item El administrador selecciona uno o varios documentos.
			\item El sistema valida los documentos y los almacena.
			\item El sistema registra la carga.
		\end{enumerate}\\
		\textbf{Postcondición}        & Los documentos están disponibles para procesado. \\
		\textbf{Excepciones}          & 
		\begin{enumerate}
			\def\labelenumi{\arabic{enumi}.}
			\tightlist
			\item Documento corrupto o sin texto legible $\Rightarrow$ Se marca como fallido y se continua con el siguiente.
		\end{enumerate}\\
		\textbf{Importancia}          & Alta.  \\
		\bottomrule
	\end{tabularx}
	\caption{\textbf{CU-9} Cargar documentos de licitaciones.}
	\label{req-09}
\end{table}

\begin{table}[p]
	\centering
	\begin{tabularx}{\linewidth}{ p{0.21\columnwidth} p{0.71\columnwidth} }
		\toprule
		\textbf{\textbf{CU-10}}    & \textbf{Importar licitaciones automáticamente}\\
		\toprule
		\textbf{Versión}              & 1.0    \\
		\textbf{Autor}                & Lydia Blanco Ruiz \\
		\textbf{Requisitos asociados} & \refreq{RF-12}, \refreq{RF-12.1}, \refreq{RF-12.2}, \refreq{RNF-3}, \refreq{RNF-3.1}\\
		\textbf{Descripción}          & Permite al administrador importar de manera automática las licitaciones a través de \eng{web scraping}. \\
		\textbf{Precondición}         & Que el administrador esté autentificado y haya un programa de \eng{web scraping} funcional. \\
		\textbf{Acciones}             &
		\begin{enumerate}
			\def\labelenumi{\arabic{enumi}.}
			\tightlist
			\item El administrador accede a <<Administrar documentos>>.
			\item El administrador pulsa el botón <<Web scraping>>.
			\item El sistema lanza el proceso de \eng{web scraping}.
			\item El sistema guarda los documentos obtenidos y registra sus metadatos.
			\item El sistema muestra un resumen con el  número de documentos descargados, los duplicados y los fallidos.
		\end{enumerate}\\
		\textbf{Postcondición}        & Los documentos importados están disponibles para procesado e indexado. \\
		\textbf{Excepciones}          & 
		\begin{enumerate}
			\def\labelenumi{\arabic{enumi}.}
			\tightlist
			\item Fallo de red $\Rightarrow$ Se registra y muestra al administrador.
			\item Documento corrupto o sin texto legible $\Rightarrow$ Se marca como fallido y se continua con el siguiente.
			\item Documento duplicado $\Rightarrow$ Se omite y se registra como duplicado.
		\end{enumerate}\\
		\textbf{Importancia}          & Media.  \\
		\bottomrule
	\end{tabularx}
	\caption{\textbf{CU-10} Importar licitaciones automáticamente.}
	\label{req-10}
\end{table}

\begin{table}[p]
	\centering
	\begin{tabularx}{\linewidth}{ p{0.21\columnwidth} p{0.71\columnwidth} }
		\toprule
		\textbf{\textbf{CU-11}}    & \textbf{Convertir los documentos a \eng{Markdown}}\\
		\toprule
		\textbf{Versión}              & 1.0    \\
		\textbf{Autor}                & Lydia Blanco Ruiz \\
		\textbf{Requisitos asociados} & \refreq{RF-2}, \refreq{RF-2.1}, \refreq{RF-2.2}, \refreq{RF-2.2.1}, \refreq{RF-2.3}, \refreq{RF-2.3.1}, \refreq{RF-2.3.2} \\
		\textbf{Descripción}          & El sistema prepara automáticamente los documentos para que puedan ser procesados por el sistema. \\
		\textbf{Precondición}         & Hay documentos cargados por el administrador. \\
		\textbf{Acciones}             &
		\begin{enumerate}
			\def\labelenumi{\arabic{enumi}.}
			\tightlist
			\item El sistema extrae texto descartando páginas y documentos sin texto extraíble.
			\item Convierte el texto a \eng{Markdown} homogéneo.
			\item Intenta conservar la estructura básica del documento original.
			\item Limpia caracteres problemáticos y normaliza espacios y saltos de línea.
		\end{enumerate}\\
		\textbf{Postcondición}        & Los documentos convertidos a \eng{Markdown} listos para el \eng{chunking}. \\
		\textbf{Excepciones}          & 
		\begin{enumerate}
			\def\labelenumi{\arabic{enumi}.}
			\tightlist
			\item Documento sin texto $\Rightarrow$ Se marca como no procesable y se continua como los demás.
		\end{enumerate}\\
		\textbf{Importancia}          & Alta.  \\
		\bottomrule
	\end{tabularx}
	\caption{\textbf{CU-11} Convertir los documentos a \eng{Markdown}.}
	\label{req-11}
\end{table}

\begin{table}[p]
	\centering
	\begin{tabularx}{\linewidth}{ p{0.21\columnwidth} p{0.71\columnwidth} }
		\toprule
		\textbf{\textbf{CU-12}}    & \textbf{Indexar documento}\\
		\toprule
		\textbf{Versión}              & 1.0    \\
		\textbf{Autor}                & Lydia Blanco Ruiz \\
		\textbf{Requisitos asociados} & \refreq{RF-3}, \refreq{RF-4}, \refreq{RF-5}, \refreq{RF-5.1}, \refreq{RF-5.2}, \refreq{RF-5.3} \\
		\textbf{Descripción}          & Permite al administrador actualizar la base de datos
vectorial cuando hay documentos nuevos. El sistema convierte los documentos procesados en \eng{chunks}, genera los \eng{embeddings} y los guarda en la base de datos vectorial con metadatos. \\
		\textbf{Precondición}         & Existen documentos procesados, una base de datos vectorial y un administrador autentificado. \\
		\textbf{Acciones}             &
		\begin{enumerate}
			\def\labelenumi{\arabic{enumi}.}
			\tightlist
			\item El administrador accede a <<Administrar documentos>>.
			\item El administrador pulsa el botón <<Actualizar Base de Datos Vectorial>>.
			\item El sistema segmenta los archivos en \eng{chunks}.
			\item Aplica solapamiento a dichos \eng{chunks}.
			\item Añade una referencia al documento así como el segmento que contiene.
			\item Genera \eng{embeddings} de cada \eng{chunk}.
			\item Lo almacena en la base de datos vectorial.
		\end{enumerate}\\
		\textbf{Postcondición}        & La colección vectorial con los nuevos documentos incluidos. \\
		\textbf{Excepciones}          & 
		\begin{enumerate}
			\def\labelenumi{\arabic{enumi}.}
			\tightlist
			\item Fallo en conseguir los \eng{embedding} de un \eng{chunk} $\Rightarrow$ Se omite el \eng{chunk} y se continúa.
			\item Error de conexión con la base de datos vectorial $\Rightarrow$ Se aborta la operación y se avisa al administrador
a través de un mensaje.
		\end{enumerate}\\
		\textbf{Importancia}          & Alta.  \\
		\bottomrule
	\end{tabularx}
	\caption{\textbf{CU-12} Indexar documento.}
	\label{req-12}
\end{table}

\begin{table}[p]
	\centering
	\begin{tabularx}{\linewidth}{ p{0.21\columnwidth} p{0.71\columnwidth} }
		\toprule
		\textbf{\textbf{CU-13}}    & \textbf{Listar documentos}\\
		\toprule
		\textbf{Versión}              & 1.0    \\
		\textbf{Autor}                & Lydia Blanco Ruiz \\
		\textbf{Requisitos asociados} & \refreq{RF-1.2}, \refreq{RNF-4} \\
		\textbf{Descripción}          & Permite al administrador ver el listado de documentos junto con los metadatos relevantes. \\
		\textbf{Precondición}         & Hay un administrador autentificado y existe al menos un documento cargado en el sistema. \\
		\textbf{Acciones}             &
		\begin{enumerate}
			\def\labelenumi{\arabic{enumi}.}
			\tightlist
			\item El administrador accede a <<Administrar documentos>>.
			\item El sistema lista los documentos con sus metadatos (nombre, tamaño, fecha de modificación y \eng{chunks}).
		\end{enumerate}\\
		\textbf{Postcondición}        & Documentos visibles para el administrador. \\
		\textbf{Excepciones}          & 
		\begin{enumerate}
			\def\labelenumi{\arabic{enumi}.}
			\tightlist
			\item No hay documentos $\Rightarrow$ Muestra un mensaje informativo <<Aún no hay documentos>>.
			\item Error de base de datos $\Rightarrow$ Muestra un error.
		\end{enumerate}\\
		\textbf{Importancia}          & Media.  \\
		\bottomrule
	\end{tabularx}
	\caption{\textbf{CU-13} Listar documentos.}
	\label{req-13}
\end{table}

\begin{table}[p]
	\centering
	\begin{tabularx}{\linewidth}{ p{0.21\columnwidth} p{0.71\columnwidth} }
		\toprule
		\textbf{\textbf{CU-14}}    & \textbf{Consultar estado del documento}\\
		\toprule
		\textbf{Versión}              & 1.0    \\
		\textbf{Autor}                & Lydia Blanco Ruiz \\
		\textbf{Requisitos asociados} & \refreq{RF-1.2}, \refreq{RNF-4} \\
		\textbf{Descripción}          & Permite al administrador ver el estado de los documentos cargados en el sistema (cargado, indexado o fallido). \\
		\textbf{Precondición}         & Hay un administrador autentificado y existe al menos un documento cargado en el sistema. \\
		\textbf{Acciones}             &
		\begin{enumerate}
			\def\labelenumi{\arabic{enumi}.}
			\tightlist
			\item El administrador accede a <<Administrar documentos>>.
			\item El sistema lista los documentos con sus metadatos (nombre, tamaño, fecha de modificación y \eng{chunks}).
			\item El sistema muestra el estado de los documentos listados.
		\end{enumerate}\\
		\textbf{Postcondición}        & Estados de los documentos visibles para el administrador. \\
		\textbf{Excepciones}          & 
		\begin{enumerate}
			\def\labelenumi{\arabic{enumi}.}
			\tightlist
			\item No hay documentos $\Rightarrow$ Muestra un mensaje informativo <<Aún no hay documentos>>.
			\item Error de base de datos $\Rightarrow$ Muestra un error.
		\end{enumerate}\\
		\textbf{Importancia}          & Media.  \\
		\bottomrule
	\end{tabularx}
	\caption{\textbf{CU-14} Consultar estado del documento.}
	\label{req-14}
\end{table}

\begin{table}[p]
	\centering
	\begin{tabularx}{\linewidth}{ p{0.21\columnwidth} p{0.71\columnwidth} }
		\toprule
		\textbf{\textbf{CU-15}}    & \textbf{Eliminar documento}\\
		\toprule
		\textbf{Versión}              & 1.0    \\
		\textbf{Autor}                & Lydia Blanco Ruiz \\
		\textbf{Requisitos asociados} & \refreq{RF-1.3}, \refreq{RF-5.5}, \refreq{RNF-5.2} \\
		\textbf{Descripción}          & Permite al administrador eliminar un documento del sistema, así como sus \eng{chunks} y \eng{embeddings} asociados en la base de datos vectorial. \\
		\textbf{Precondición}         & Hay un administrador autentificado y existe al menos un documento cargado en el sistema. \\
		\textbf{Acciones}             &
		\begin{enumerate}
			\def\labelenumi{\arabic{enumi}.}
			\tightlist
			\item El administrador accede a <<Administrar documentos>>.
			\item El sistema lista los documentos con sus metadatos (nombre, tamaño, fecha de modificación y \eng{chunks}) y estado.
			\item El administrador pulsa el botón <<Borrar>>.
			\item El sistema pide confirmación.
			\item El administrador confirma el borrado.
			\item El sistema borra el documento y sus metadatos. Además, si el documento esta indexado borra sus \eng{chunks} y sus \eng{embeddings}.
			\item El sistema actualiza el listado.
		\end{enumerate}\\
		\textbf{Postcondición}        & Documento eliminado y sus \eng{chunks} y \eng{embeddings} borrados de la base de datos vectorial. \\
		\textbf{Excepciones}          & 
		\begin{enumerate}
			\def\labelenumi{\arabic{enumi}.}
			\tightlist
			\item El documento no existe $\Rightarrow$ Se muestra un aviso. 
			\item Error al borra $\Rightarrow$ No se borra y se informa al administrador.
			\item Error al borrar en la base de datos vectorial $\Rightarrow$ No se borra el documento y se informa al administrador.
		\end{enumerate}\\
		\textbf{Importancia}          & Alta.  \\
		\bottomrule
	\end{tabularx}
	\caption{\textbf{CU-15} Eliminar documento.}
	\label{req-15}
\end{table}

\begin{table}[p]
	\centering
	\begin{tabularx}{\linewidth}{ p{0.21\columnwidth} p{0.71\columnwidth} }
		\toprule
		\textbf{\textbf{CU-16}}    & \textbf{Descargar documento}\\
		\toprule
		\textbf{Versión}              & 1.0    \\
		\textbf{Autor}                & Lydia Blanco Ruiz \\
		\textbf{Requisitos asociados} & \refreq{RF-1}, \refreq{RF-1.4}, \refreq{RNF-5.2} \\
		\textbf{Descripción}          & Permite al administrador descargar un documento del sistema. \\
		\textbf{Precondición}         & Hay un administrador autentificado y existe al menos un documento cargado en el sistema. \\
		\textbf{Acciones}             &
		\begin{enumerate}
			\def\labelenumi{\arabic{enumi}.}
			\tightlist
			\item El administrador accede a <<Administrar documentos>>.
			\item El sistema lista los documentos con sus metadatos (nombre, tamaño, fecha de modificación y \eng{chunks}) y estado.
			\item El administrador pulsa el botón <<Descargar>>.
			\item Se descarga el documento.
		\end{enumerate}\\
		\textbf{Postcondición}        & Documento descargado en el equipo del administrador. \\
		\textbf{Excepciones}          & 
		\begin{enumerate}
			\def\labelenumi{\arabic{enumi}.}
			\tightlist
			\item El documento no existe $\Rightarrow$ Se muestra un aviso. 
			\item Error al descargar $\Rightarrow$ Se informa al administrador.
		\end{enumerate}\\
		\textbf{Importancia}          & Media.  \\
		\bottomrule
	\end{tabularx}
	\caption{\textbf{CU-16} Descargar documento.}
	\label{req-16}
\end{table}

\begin{table}[p]
	\centering
	\begin{tabularx}{\linewidth}{ p{0.21\columnwidth} p{0.71\columnwidth} }
		\toprule
		\textbf{\textbf{CU-17}}    & \textbf{ Gestión de usuarios}\\
		\toprule
		\textbf{Versión}              & 1.0    \\
		\textbf{Autor}                & Lydia Blanco Ruiz \\
		\textbf{Requisitos asociados} & \refreq{RF-10.3}, \refreq{RF-10.4} \\
		\textbf{Descripción}          & El administrador debe poder gestionar a los usuarios de la aplicación y sus permisos. \\
		\textbf{Precondición}         & Hay un administrador autentificado. \\
		\textbf{Acciones}             &
		\begin{enumerate}
			\def\labelenumi{\arabic{enumi}.}
			\tightlist
			\item El administrador accede a la página de <<Iniciar sesión>>.
			\item El administrador introduce su \eng{email} y contraseña.
			\item El sistema le autentifica como administrador.
			\item El administrador accede a <<Administrar usuarios>>.
		\end{enumerate}\\
		\textbf{Postcondición}        & Listado de usuarios para editar. \\
		\textbf{Excepciones}          & 
		\begin{enumerate}
			\def\labelenumi{\arabic{enumi}.}
			\tightlist
			\item Intento de acceso sin ser administrador $\Rightarrow$ Se deniega.
			\item Si no hay usuarios registrados en la aplicación $\Rightarrow$ Mostrar mensaje de que aún no hay usuarios.
		\end{enumerate}\\
		\textbf{Importancia}          & Media.  \\
		\bottomrule
	\end{tabularx}
	\caption{\textbf{CU-17} Gestión de usuarios.}
	\label{req-17}
\end{table}

\begin{table}[p]
	\centering
	\begin{tabularx}{\linewidth}{ p{0.21\columnwidth} p{0.71\columnwidth} }
		\toprule
		\textbf{\textbf{CU-18}}    & \textbf{ Crear usuario}\\
		\toprule
		\textbf{Versión}              & 1.0    \\
		\textbf{Autor}                & Lydia Blanco Ruiz \\
		\textbf{Requisitos asociados} & \refreq{RF-10.3}, \refreq{RF-10.4} \\
		\textbf{Descripción}          & El administrador crea usuarios desde la web. \\
		\textbf{Precondición}         & Hay un administrador autentificado. \\
		\textbf{Acciones}             &
		\begin{enumerate}
			\def\labelenumi{\arabic{enumi}.}
			\tightlist
			\item El administrador accede a <<Administrar usuarios>>.
			\item El administrador pulsa el botón <<Añadir usuario>>.
			\item El administrador introduce el nombre, \eng{email}, contraseña y si es administrador o no.
			\item El sistema valida el formato y que el \eng{email} no esté duplicado.
			\item El sistema almacena el usuario con la contraseña encriptada.
			\item El sistema redirige a la página de <<Gestión de usuarios>>.
		\end{enumerate}\\
		\textbf{Postcondición}        & Usuario creado. \\
		\textbf{Excepciones}          & 
		\begin{enumerate}
			\def\labelenumi{\arabic{enumi}.}
			\tightlist
			\item Intento de acceso sin ser administrador $\Rightarrow$ Se deniega.
			\item Tiene algún campo vacío $\Rightarrow$ Muestra un error con el mensaje <<Completa este campo>>.
			\item El \eng{email} ya existe $\Rightarrow$ Muestra un error con el mensaje <<Ya existe un usuario con ese \eng{email}>>.
			\item El \eng{email} no cumple la validación $\Rightarrow$ Muestra un error con el mensaje <<Invalid \eng{email} address>>.
			\item La contraseña no cumple la política $\Rightarrow$ Muestra un error con el mensaje <<Aumenta la longitud del texto a 6 caracteres como mínimo>>.
			\item El nombre tiene menos de 2 caracteres $\Rightarrow$ Muestra un error con el mensaje <<Aumenta la longitud del texto a 2 caracteres o más>>.
		\end{enumerate}\\
		\textbf{Importancia}          & Media.  \\
		\bottomrule
	\end{tabularx}
	\caption{\textbf{CU-18} Crear usuario.}
	\label{req-18}
\end{table}

\begin{table}[p]
	\centering
	\begin{tabularx}{\linewidth}{ p{0.21\columnwidth} p{0.71\columnwidth} }
		\toprule
		\textbf{\textbf{CU-19}}    & \textbf{ Borrar usuario}\\
		\toprule
		\textbf{Versión}              & 1.0    \\
		\textbf{Autor}                & Lydia Blanco Ruiz \\
		\textbf{Requisitos asociados} & \refreq{RF-10.3}, \refreq{RF-10.4} \\
		\textbf{Descripción}          & El administrador borra usuarios desde la web. \\
		\textbf{Precondición}         & Hay un administrador autentificado y usuarios previamente registrados en la aplicación. \\
		\textbf{Acciones}             &
		\begin{enumerate}
			\def\labelenumi{\arabic{enumi}.}
			\tightlist
			\item El administrador accede a <<Administrar usuarios>>.
			\item El administrador pulsa el botón <<Borrar>> del usuario que desee eliminar.
			\item El sistema pide confirmación antes de borrar el usuario.
			\item El administrador pulsa el botón <<Aceptar>>.
			\item El sistema borra al usuario y sus datos de la base de datos.
		\end{enumerate}\\
		\textbf{Postcondición}        & Usuario borrado. \\
		\textbf{Excepciones}          & 
		\begin{enumerate}
			\def\labelenumi{\arabic{enumi}.}
			\tightlist
			\item Fallo al conectar con la base de datos $\Rightarrow$ No se borra el usuario y se informa del problema al administrador.
		\end{enumerate}\\
		\textbf{Importancia}          & Media.  \\
		\bottomrule
	\end{tabularx}
	\caption{\textbf{CU-19} Borrar usuario.}
	\label{req-19}
\end{table}

\begin{table}[p]
	\centering
	\begin{tabularx}{\linewidth}{ p{0.21\columnwidth} p{0.71\columnwidth} }
		\toprule
		\textbf{\textbf{CU-20}}    & \textbf{ Modificar rol del usuario}\\
		\toprule
		\textbf{Versión}              & 1.0    \\
		\textbf{Autor}                & Lydia Blanco Ruiz \\
		\textbf{Requisitos asociados} & \refreq{RF-10.3}, \refreq{RF-10.4} \\
		\textbf{Descripción}          & El administrador puede cambiar el rol de los usuarios desde la web. \\
		\textbf{Precondición}         & Hay un administrador autentificado y usuarios previamente registrados en la aplicación. \\
		\textbf{Acciones}             &
		\begin{enumerate}
			\def\labelenumi{\arabic{enumi}.}
			\tightlist
			\item El administrador accede a <<Administrar usuarios>>.
			\item El administrador pulsa el botón <<Hacer admin>>/<<Quitar admin>> del usuario que desee cambiar de rol.
			\item El sistema cambia el rol del usuario.
		\end{enumerate}\\
		\textbf{Postcondición}        & Usuario con el rol cambiado. \\
		\textbf{Excepciones}          & 
		\begin{enumerate}
			\def\labelenumi{\arabic{enumi}.}
			\tightlist
			\item Fallo al conectar con la base de datos $\Rightarrow$ No se cambia el rol del usuario y se informa del problema al administrador.
		\end{enumerate}\\
		\textbf{Importancia}          & Baja.  \\
		\bottomrule
	\end{tabularx}
	\caption{\textbf{CU-20} Modificar rol del usuario.}
	\label{req-20}
\end{table}

\begin{table}[p]
	\centering
	\begin{tabularx}{\linewidth}{ p{0.21\columnwidth} p{0.71\columnwidth} }
		\toprule
		\textbf{\textbf{CU-21}}    & \textbf{ Editar usuario}\\
		\toprule
		\textbf{Versión}              & 1.0    \\
		\textbf{Autor}                & Lydia Blanco Ruiz \\
		\textbf{Requisitos asociados} & \refreq{RF-10.5}, \refreq{RNF-5.1}, \refreq{RNF-5.2} \\
		\textbf{Descripción}          & Permite que un usuario edite sus datos y contraseña. \\
		\textbf{Precondición}         & Hay un usuario autentificado. \\
		\textbf{Acciones}             &
		\begin{enumerate}
			\def\labelenumi{\arabic{enumi}.}
			\tightlist
			\item El usuario accede a <<Editar usuario>>.
			\item Modifica los campos que desee (nombre, \eng{email} o contraseña).
			\item El sistema valida los formatos y reglas.
			\item El sistema guarda los cambios. Si hay cambio de contraseña la almacena encriptada.
		\end{enumerate}\\
		\textbf{Postcondición}        & Datos del usuario actualizados. \\
		\textbf{Excepciones}          & 
		\begin{enumerate}
			\def\labelenumi{\arabic{enumi}.}
			\tightlist
			\item Se intenta cambiar el \eng{email} por otro que ya existe en el sistema $\Rightarrow$ Muestra un error con el mensaje <<Ya existe un usuario con ese \eng{email}>>.
			\item El \eng{email} no cumple la validación $\Rightarrow$ Muestra un error con el mensaje <<Invalid \eng{email} address>>.
			\item La contraseña no cumple la política $\Rightarrow$ Muestra un error con el mensaje <<Aumenta la longitud del texto a 6 caracteres como mínimo>>.
			\item El nombre tiene menos de 2 caracteres $\Rightarrow$ Muestra un error con el mensaje <<Aumenta la longitud del texto a 2 caracteres o más>>.
		\end{enumerate}\\
		\textbf{Importancia}          & Baja.  \\
		\bottomrule
	\end{tabularx}
	\caption{\textbf{CU-21} Editar usuario.}
	\label{req-21}
\end{table}


\begin{table}[p]
	\centering
	\begin{tabularx}{\linewidth}{ p{0.21\columnwidth} p{0.71\columnwidth} }
		\toprule
		\textbf{\textbf{CU-22}}    & \textbf{ Realizar consulta guiada}\\
		\toprule
		\textbf{Versión}              & 1.0    \\
		\textbf{Autor}                & Lydia Blanco Ruiz \\
		\textbf{Requisitos asociados} & \refreq{RF-11}, \refreq{RF-11.3}, \refreq{RF-6}, \refreq{RF-7}, \refreq{RF-8}, \refreq{RF-9} \\
		\textbf{Descripción}          & Permite al usuario realizar consultas utilizando un formulario que llama a \eng{prompts} especializados. \\
		\textbf{Precondición}         & Hay un usuario autentificado y existen documentos indexados en el sistema. \\
		\textbf{Acciones}             &
		\begin{enumerate}
			\def\labelenumi{\arabic{enumi}.}
			\tightlist
			\item El usuario accede a la interfaz de consulta.
			\item El usuario selecciona una pregunta del formulario (resumen, importes, solvencia, garantías, criterios de adjudicación, etc.).
			\item El usuario pulsa el botón <<Preguntar>>.
			\item El sistema genera el \eng{embedding} de la consulta.
			\item El sistema recupera los \eng{chunks} más relevantes.
			\item El sistema construye el \eng{prompt} aumentado según el perfil seleccionado.
			\item El sistema envía el \eng{prompt} al \eng{LLM}.
			\item El sistema muestra la respuesta generada junto con la trazabilidad.
		\end{enumerate}\\
		\textbf{Postcondición}        & El usuario ve la respuesta y se almacena en el historial. \\
		\textbf{Excepciones}          & 
		\begin{enumerate}
			\def\labelenumi{\arabic{enumi}.}
			\tightlist
			\item No existen documentos indexados $\Rightarrow$ Muestra un mensaje de error.
			\item La consulta está vacía $\Rightarrow$ Muestra un mensaje indicando que debe introducirse una consulta.
			\item No se encuentra contexto relevante $\Rightarrow$ El sistema informa de que no dispone de información suficiente.
			\item El \eng{LLM} no responde o se produce un \eng{timeout} $\Rightarrow$ Muestra un mensaje de error controlado.
		\end{enumerate}\\
		\textbf{Importancia}          & Alta.  \\
		\bottomrule
	\end{tabularx}
	\caption{\textbf{CU-22} Realizar consulta guiada.}
	\label{req-22}
\end{table}

\begin{table}[p]
	\centering
	\begin{tabularx}{\linewidth}{ p{0.21\columnwidth} p{0.71\columnwidth} }
		\toprule
		\textbf{\textbf{CU-23}}    & \textbf{ Cambiar idioma}\\
		\toprule
		\textbf{Versión}              & 1.0    \\
		\textbf{Autor}                & Lydia Blanco Ruiz \\
		\textbf{Requisitos asociados} & \refreq{RNF-16} \\
		\textbf{Descripción}          & Permite al usuario cambiar el idioma de la interfaz de la aplicación. \\
		\textbf{Precondición}         & Un usuario accede a la aplicación \eng{web}. \\
		\textbf{Acciones}             &
		\begin{enumerate}
			\def\labelenumi{\arabic{enumi}.}
			\tightlist
			\item El usuario accede a la aplicación \eng{web}.
			\item El usuario selecciona el idioma deseado
			\item El sistema recarga la interfaz mostrando los textos traducidos.
			\item Si el usuario está autenticado el sistema guarda la preferencia de idioma.
		\end{enumerate}\\
		\textbf{Postcondición}        & La aplicación se muestra en el idioma seleccionado. \\
		\textbf{Excepciones}          & 
		\begin{enumerate}
			\def\labelenumi{\arabic{enumi}.}
			\tightlist
			\item El idioma seleccionado no está disponible $\Rightarrow$ El sistema mantiene el idioma actual.
			\item Error al guardar la preferencia $\Rightarrow$ Muestra un mensaje de error.
		\end{enumerate}\\
		\textbf{Importancia}          & Media.  \\
		\bottomrule
	\end{tabularx}
	\caption{\textbf{CU-23} Cambiar idioma.}
	\label{req-23}
\end{table}

\begin{table}[p]
	\centering
	\begin{tabularx}{\linewidth}{ p{0.21\columnwidth} p{0.71\columnwidth} }
		\toprule
		\textbf{\textbf{CU-24}}    & \textbf{ Cambiar tema visual}\\
		\toprule
		\textbf{Versión}              & 1.0    \\
		\textbf{Autor}                & Lydia Blanco Ruiz \\
		\textbf{Requisitos asociados} & \refreq{RNF-6}, \refreq{RNF-17}\\
		\textbf{Descripción}          & Permite al usuario cambiar entre el modo claro y el modo oscuro. Por defecto se muestra en el modo de l sistema. \\
		\textbf{Precondición}         & Un usuario accede a la aplicación \eng{web}. \\
		\textbf{Acciones}             &
		\begin{enumerate}
			\def\labelenumi{\arabic{enumi}.}
			\tightlist
			\item El usuario accede a la aplicación \eng{web}.
			\item El usuario selecciona el tema deseado.
			\item El sistema recarga la interfaz mostrando los elementos visuales del tema seleccionado.
			\item Si el usuario está autenticado el sistema guarda la preferencia de modo.
		\end{enumerate}\\
		\textbf{Postcondición}        & La aplicación se muestra utilizando el tema seleccionado (claro/oscuro). \\
		\textbf{Excepciones}          & 
		\begin{enumerate}
			\def\labelenumi{\arabic{enumi}.}
			\tightlist
			\item Error al guardar la preferencia $\Rightarrow$ Muestra un mensaje de error.
		\end{enumerate}\\
		\textbf{Importancia}          & Media.  \\
		\bottomrule
	\end{tabularx}
	\caption{\textbf{CU-24} Cambiar tema visual.}
	\label{req-24}
\end{table}

\begin{table}[p]
	\centering
	\begin{tabularx}{\linewidth}{ p{0.21\columnwidth} p{0.71\columnwidth} }
		\toprule
		\textbf{\textbf{CU-25}}    & \textbf{ Visualizar estadísticas}\\
		\toprule
		\textbf{Versión}              & 1.0    \\
		\textbf{Autor}                & Lydia Blanco Ruiz \\
		\textbf{Requisitos asociados} & \refreq{RF-13}, \refreq{RNF-6}\\
		\textbf{Descripción}          & Permite visualizar estadísticas y gráficas relacionadas con el uso de la aplicación por parte de los usuarios y comparativas entre los modelos \eng{LLM} cargados en el sistema. \\
		\textbf{Precondición}         & Hay un usuario autentificado y existen datos registrados en el sistema. \\
		\textbf{Acciones}             &
		\begin{enumerate}
			\def\labelenumi{\arabic{enumi}.}
			\tightlist
			\item El usuario accede a la página de estadísticas.
			\item El sistema recupera los datos almacenados.
			\item El sistema genera gráficas y métricas mediante \file{D3.js}.
			\item El usuario ve las gráficas disponibles para su rol.
		\end{enumerate}\\
		\textbf{Postcondición}        & La aplicación muestra las gráficas y estadísticas. \\
		\textbf{Excepciones}          & 
		\begin{enumerate}
			\def\labelenumi{\arabic{enumi}.}
			\tightlist
			\item No existen datos suficientes $\Rightarrow$ El sistema muestra un mensaje informativo.
			\item Error al cargar las estadísticas $\Rightarrow$ Muestra un mensaje de error.
		\end{enumerate}\\
		\textbf{Importancia}          & Media.  \\
		\bottomrule
	\end{tabularx}
	\caption{\textbf{CU-25} Visualizar estadísticas.}
	\label{req-25}
\end{table}

\begin{table}[p]
	\centering
	\begin{tabularx}{\linewidth}{ p{0.21\columnwidth} p{0.71\columnwidth} }
		\toprule
		\textbf{\textbf{CU-26}}    & \textbf{ Ejecutar evaluación \eng{RAG}}\\
		\toprule
		\textbf{Versión}              & 1.0    \\
		\textbf{Autor}                & Lydia Blanco Ruiz \\
		\textbf{Requisitos asociados} & \refreq{RF-14}, \refreq{RF-14.1}, \refreq{RNF-11} \\
		\textbf{Descripción}          & Permite al administrador ejecutar evaluaciones automáticas sobre la calidad del \eng{RAG}. \\
		\textbf{Precondición}         & Hay un administrador autentificado y existen documentos indexados en la base de datos vectorial. \\
		\textbf{Acciones}             &
		\begin{enumerate}
			\def\labelenumi{\arabic{enumi}.}
			\tightlist
			\item El administrador accede a la página de evaluación.
			\item El administrador selecciona el modelo y configuración de evaluación.
			\item Se ejecuta el \eng{pipeline} de evaluación.		
			\item El sistema almacena los resultados obtenidos.
			\item El sistema muestra las métricas y resultados de evaluación.
			\end{enumerate}\\
		\textbf{Postcondición}        & Los resultados de la evaluación quedan almacenados y disponibles para su consulta. \\
		\textbf{Excepciones}          & 
		\begin{enumerate}
			\def\labelenumi{\arabic{enumi}.}
			\tightlist
			\item Error durante la ejecución de las métricas $\Rightarrow$ El sistema registra el error y detiene la evaluación.
			\item El modelo \eng{LLM} no responde $\Rightarrow$ Muestra un mensaje de \eng{timeout} o error controlado.
		\end{enumerate}\\
		\textbf{Importancia}          & Alta.  \\
		\bottomrule
	\end{tabularx}
	\caption{\textbf{CU-26} Ejecutar evaluación \eng{RAG}.}
	\label{req-26}
\end{table}